% !TeX root = ../kaehler-einstein-degeneration.tex
% Current content: 2026-11-23, v1.2.  Explicit real-Hermitian ellipticity constants for Savin H2.
\section{三重点の開正方形領域における actual KE 比較}
\label{sec:moving-triple-comparison}

本節では，\cref{ass:Gd} を満たす固定次数 $d\ge5$ の全平面退化
$X_t^{(d)}$ を考え，一つの logarithmic depth が他の二つより大きい三重点領域を扱う．
$\omega_t$ と $g_t$ は $\Ric(\omega_t)=-\omega_t$ と規格化した actual KE form と実計量とする．
以下の定数は固定した $d$ と退化族に依存してよく，次数を動かす一様性は主張しない．
この領域では三重点 chart を出ない expanding product annuli を直接構成できる．
大域 volume 下界を境界値として用いる負 KE 下 barrier と，complete annulus 上 barrier により，
volume ratio を $1$ に挟み込む．その結果として得られる\emph{小さい} potential に
Savin の定理を適用し，actual metric の tensor 比較を導く．
単に有界な volume ratio から tensor 比較を推論する議論とは異なり，
\cref{coarse:prop:density-counterexample} の反例と矛盾しない．
本節は exterior の tensor-comparison 仮定や $\mathrm{(PBE)}_d$，
任意半径の complete-envelope package を用いない．
依存する大域入力は \cref{coarse:cor:degree-d-volume-lower} である．
五次では \cref{coarse:cor:global-volume-lower} を用いても同じ証明になる．
以下の組合せによる結果の新規性は未確認である．

\subsection{正確な product domain と粗い volume 下界}

Source を含まない三重点で，正則座標を取り直して
\begin{equation}\label{moving:eq:exact-triple}
 \zeta_0\zeta_1\zeta_2=t,\qquad
 \ell_i=-\log|\zeta_i|^2,\qquad
 \ell_0+\ell_1+\ell_2=L_t
\end{equation}
とする．この座標変換に現れる units と canonical residue frame の係数は，
固定した小さい三重点 chart 上で一様に正の上下界を持つ．
以下では $s_1=\ell_1$, $s_2=\ell_2$, $\theta_i=\arg \zeta_i$ と書き，
独立な二つの punctured-disc cusp の積を
\begin{equation}\label{moving:eq:product-cusp}
 h_{\mathrm{prod}}
 =\frac{ds_1^2+4d\theta_1^2}{s_1^2}
  +\frac{ds_2^2+4d\theta_2^2}{s_2^2}
\end{equation}
とする．同じ記号は対応する K\"ahler form にも用いる．
これは \eqref{coarse:eq:moving-cusp-product} の $\omega_t^{\mathrm{prod}}$ と
$g_t^{\mathrm{prod}}$ であり，三因子を制約した $g_t^{\mathrm{cusp}}$ とは区別する．
前節の $dV_{\log}$ を用いると，正規化は
\[
 h_{\mathrm{prod}}^2=\frac{8}{s_1^2s_2^2}dV_{\log},
 \qquad
 \frac{\mathcal V_t}{h_{\mathrm{prod}}^2}=\frac{L_t^2}{\ell_0^2}
\]
である．各因子の Gaussian curvature は $-1$ であり，
$\Ric(h_{\mathrm{prod}})=-h_{\mathrm{prod}}$ である．

中心深度 $a,b>0$ と $\rho>0$ に対し，全 angular orbits を含む domain を
\begin{equation}\label{moving:eq:rectangle}
 D_t(a,b;\rho)
 =\{e^{-\rho}a<s_1<e^\rho a,\quad
      e^{-\rho}b<s_2<e^\rho b\}
\end{equation}
と置く．$\zeta_0=t/(\zeta_1\zeta_2)$ は従属座標である．
固定した十分大きい $m_0$ に対して
\begin{equation}\label{moving:eq:rectangle-protection}
 e^{-\rho}\min(a,b)\ge m_0,
 \qquad e^\rho(a+b)\le \frac{L_t}{2}
\end{equation}
ならば，$\overline{D_t(a,b;\rho)}$ は同じ三重点 chart 内にある．
実際，$s_1,s_2\ge m_0$ と $\ell_0=L_t-s_1-s_2\ge L_t/2$ により，
三つの $|\zeta_i|$ がすべて chart の半径より小さい．
この包含は全位相について成立する．

\begin{lemma}[Expanding rectangle 上の volume 下界]
\label{moving:lem:coarse-volume-lower}
固定した $m_0$ を十分大きく取ると，
\eqref{moving:eq:rectangle-protection} を満たす全 $t,a,b,\rho$ に対し
\begin{equation}\label{moving:eq:coarse-volume-lower}
 \omega_t^2\ge c_* h_{\mathrm{prod}}^2
 \qquad\text{on }\overline{D_t(a,b;\rho)}
\end{equation}
が成り立つ．$c_*>0$ は $t,a,b,\rho$ と angular phases に依存しない．
\end{lemma}

\begin{proof}
三重点で消える三つの標識について，固定 Hermitian norm と
\eqref{moving:eq:exact-triple} の座標との差は bounded units である．
従って capped depths は，この rectangle 上で
\[
 \alpha_0=\frac{L_t}{2}+O(1),\qquad
 \alpha_1=s_1+O(1),\qquad
 \alpha_2=s_2+O(1)
\]
を満たす．残る標識の depths は正の定数の間にある．
Canonical residue frame は，bounded nonvanishing factor を除いて
\[
 \frac{d\zeta_1}{\zeta_1}\wedge\frac{d\zeta_2}{\zeta_2}
\]
である．従って $\Omega_{t,d}^{\mathrm{low}}=L_t^2\prod_{i=1}^d\alpha_i^{-2}\Omega_{0,t}^{(d)}$ は
$h_{\mathrm{prod}}^2$ と一様に比較可能である．
\Cref{coarse:cor:degree-d-volume-lower} を適用すればよい．
ここで比較するのは density $\Omega_{t,d}^{\mathrm{low}}$ であり，
その curvature form の tensor 比較は用いていない．
\end{proof}

\subsection{Shifted-sinh barrier による volume ratio の挟み込み}

\begin{lemma}[二つの有限端を持つ負 KE 下 barrier]
\label{moving:lem:sinh-barrier}
$A>0$, $\rho>0$, $M>0$ とし，
\[
 s_-=e^{-\rho}A,\qquad s_+=e^\rho A,
 \qquad c=M e^{-\rho}A,\qquad
 \lambda=\frac{M e^{-\rho}}{A}
\]
と置く．有限 annulus $s_-\le s\le s_+$ 上の計量
\begin{equation}\label{moving:eq:sinh-barrier}
 h^-_{A,\rho,M}
 = f^-(s)(ds^2+4d\theta^2),\qquad
 f^-(s)=\frac{\lambda^2}{\sinh^2(\lambda(s+c))}
\end{equation}
は Gaussian curvature $-1$ を持ち，
$h^-_{A,\rho,M}\le h_{\mathrm{cusp},A}:=(ds^2+4d\theta^2)/s^2$ を満たす．
Ratio $r^-(s)=s^2f^-(s)$ は二つの端で
\begin{equation}\label{moving:eq:sinh-face-ratios}
 r^-(s_-)\le(1+M)^{-2},\qquad
 r^-(s_+)\le\left(\frac{M}{\sinh M}\right)^2
\end{equation}
を満たす．任意の固定 $r>0$ について，$\rho\to\infty$ のとき
\begin{equation}\label{moving:eq:sinh-center-ratio}
 \sup_{e^{-r}A\le s\le e^r A}|r^-(s)-1|
 \le C_{M,r}e^{-\rho}
\end{equation}
であり，定数は $A$ に依存しない．
\end{lemma}

\begin{proof}
$(\log f^-)''=2f^-$ なので，conformal metric
$f^-(s)(ds^2+4d\theta^2)$ の曲率は $-1$ である．
また $\sinh x\ge x$ より
\[
 r^-(s)
 =\left(\frac{s}{s+c}\right)^2
   \left(\frac{\lambda(s+c)}{\sinh(\lambda(s+c))}\right)^2
 \le\left(\frac{s}{s+c}\right)^2\le1.
\]
低い端では $c=M s_-$ なので最初の評価を得る．
高い端では $\lambda s_+=M$ なので第二の評価を得る．
固定 inner interval では
$c/s\le M e^{r-\rho}$ と
$\lambda(s+c)\le M e^{r-\rho}+M^2e^{-2\rho}$ である．
この式と $x/\sinh x=1+O(x^2)$ から
\eqref{moving:eq:sinh-center-ratio} が従う．
\end{proof}

同じ有限 annulus 自身の complete curvature $-1$ metric は
\begin{equation}\label{moving:eq:complete-annulus-upper}
 h^+_{A,\rho}
 = f^+(s)(ds^2+4d\theta^2),\qquad
 f^+(s)=\left(\frac{\pi}{s_+-s_-}\right)^2
 \frac{1}{\sin^2\!\left(\pi\frac{s-s_-}{s_+-s_-}\right)}.
\end{equation}
これは人工境界で発散する．$r^+(s)=s^2f^+(s)$ とすれば
\begin{equation}\label{moving:eq:complete-center-ratio}
 \sup_{e^{-r}A\le s\le e^r A}|r^+(s)-1|
 \le C_r e^{-\rho}
\end{equation}
が成立する．実際，$x=\pi(s-s_-)/(s_+-s_-)$ と置くと
\[
 r^+(s)=\left(\frac{s}{s-s_-}\right)^2
        \left(\frac{x}{\sin x}\right)^2,
 \qquad \frac{s_-}{s}=O_r(e^{-\rho}),\quad x=O_r(e^{-\rho}).
\]

\begin{proposition}[Sharp volume squeeze]
\label{coarse:lem:annular-volume-squeeze}
\eqref{moving:eq:rectangle-protection} を仮定する．
$c_*$ は \cref{moving:lem:coarse-volume-lower} の定数とし，
$M$ を固定して
\[
 \max\left\{(1+M)^{-2},
       \left(\frac{M}{\sinh M}\right)^2\right\}\le c_*
\]
を満たすように取る．Product metrics を
\[
 h^- = h^-_{a,\rho,M}\oplus h^-_{b,\rho,M},\qquad
 h^+ = h^+_{a,\rho}\oplus h^+_{b,\rho}
\]
と置くと，$D_t(a,b;\rho)$ 上で
\begin{equation}\label{moving:eq:volume-sandwich}
 (h^-)^2\le\omega_t^2\le(h^+)^2
\end{equation}
である．従って任意の固定 $r>0$ に対し，$\rho$ が十分大きいとき
\begin{equation}\label{moving:eq:volume-small-potential}
 \left\|\log\frac{\omega_t^2}{h_{\mathrm{prod}}^2}
 \right\|_{C^0(D_t(a,b;r))}
 \le C_r e^{-\rho}.
\end{equation}
全ての定数と評価は angular phases について一様である．
\end{proposition}

\begin{proof}
$h^-$ と $h^+$ はいずれも負 KE product である．
Rectangle の各 boundary face では，一方の因子の ratio が
\eqref{moving:eq:sinh-face-ratios} により $c_*$ 以下であり，
他方の ratio は $1$ 以下である．従って
\[
 (h^-)^2\le c_*h_{\mathrm{prod}}^2\le\omega_t^2
 \quad\text{on }\partial D_t(a,b;\rho).
\]
もし $u^-:=\log(\omega_t^2/(h^-)^2)$ が負の最小値を持てば，その点は内部にあり，
Ricci forms の差から
\[
 \ddc u^- =\omega_t-h^-\ge0
\]
となる．従ってその点で $\omega_t^2\ge(h^-)^2$ となり矛盾する．
これで下比較を得る．

上比較について，actual metric は rectangle の閉包の近傍で滑らかであり，
$h^+$ の density は全ての人工境界で発散するので，
$u^+:=\log(\omega_t^2/(h^+)^2)$ は境界へ近づくと $-\infty$ へ向かう．
正の内部最大点があれば
$\ddc u^+=\omega_t-h^+\le0$ より
$\omega_t^2\le(h^+)^2$ となり矛盾する．
両比較とも全位相を含む有限 annulus product 上の通常の最大値原理であり，
actual metric に関する境界までの距離を仮定していない．

Volume ratios は一変数の ratios の積である．
\eqref{moving:eq:sinh-center-ratio} と
\eqref{moving:eq:complete-center-ratio} を用いれば
\eqref{moving:eq:volume-small-potential} を得る．
\end{proof}

\subsection{Small potential から tensor 比較へ}

\begin{lemma}[Savin の定理による負 KE の局所安定性]
\label{coarse:lem:small-volume-tensor}
固定複素 ball 上の滑らかな負 KE metric $h$ を取り，
その係数は一様に正で全固定階の bounds を持つとする．
同じ ball 上の滑らかな負 KE metrics $\omega_\nu$ に対し
\[
 v_\nu=\log\frac{\omega_\nu^2}{h^2},\qquad
 \|v_\nu\|_{C^0}\longrightarrow0
\]
ならば，任意の内側 ball 上で
$\omega_\nu\to h$ が $C^\infty$ の意味で成立する．
さらに，十分小さい $\|v_\nu\|_{C^0}$ に対して
\begin{equation}\label{moving:eq:linear-smallness}
 \|\omega_\nu-h\|_{C^k(B')}
 \le C_{k,B'}\|v_\nu\|_{C^0(B)}
\end{equation}
となる．この主張は $h$ が一様な上記 bounds を持つ族を動く場合も一様である．
\end{lemma}

\begin{proof}
Ricci forms の差より
\begin{equation}\label{moving:eq:small-KE-equation}
 \ddc v_\nu=\omega_\nu-h,
 \qquad (h+\ddc v_\nu)^2=e^{v_\nu}h^2.
\end{equation}
正則座標で $h$ の Hermitian matrix を $H(x)$ とし，
real Hessian を complex Hessian へ送る線形写像を $\mathcal H$ とする．
Hermitian matrix $A$ に対して
\[
 \det_+(A)=
 \begin{cases}
  \det A,& A\ge0,\\
  0,& A\not\ge0
 \end{cases}
\]
と定義し，real symmetric matrix $N$，実数 $z$ に対して
\begin{equation}\label{moving:eq:Savin-operator}
 F(N,p,z,x)
 =\frac{\det_+(H(x)+\mathcal H(N))}{\det H(x)}-e^z
\end{equation}
と置く．$F$ は $p$ に依存しない．
$\det_+$ は全 Hermitian matrices 上で連続であり，
positive semidefinite 方向に単調である．
また $N\ge0$ なら $\mathcal H(N)\ge0$ であり，
$\tr_{\CC}\mathcal H(N)$ は $\tr_{\RR}N$ の固定正倍である．
従って $F$ は Savin の hypothesis H1 を満たす．

Hypothesis H2 の一様楕円性を実方向の定数まで確認する．
ここでは複素次元が $2$ なので，$\mathcal H$ は実対称
$4\times4$ 行列から Hermitian $2\times2$ 行列への写像である．
実 Hessian の $xx,xy,yy$ block を用いれば
\[
 \mathcal H(P)=\frac14(P_{xx}+P_{yy}
              +\sqrt{-1}(P_{xy}-P_{xy}^{\,T})),
 \quad
 \operatorname{tr}_{\CC}\mathcal H(P)=\frac14\operatorname{tr}_{\RR}P .
\]
$P\ge0$ なら $\mathcal H(P)\ge0$ であり，
$\|P\|\le\operatorname{tr}_{\RR}P\le4\|P\|$ である．
$H\ge\lambda I$ より，$N$ と $P$ を十分小さくし，
$B_s=H+\mathcal H(N+sP)$ $(0\le s\le1)$ を考えると，
$B_s$ は一様な上下の正の固有値界を持ち，
$\det B_s/\det H$ も正の上下界を持つ．従って
\[
 \frac{d}{ds}F(N+sP,p,z,x)
 =\frac{\det B_s}{\det H}
       \operatorname{tr}_{\CC}(B_s^{-1}\mathcal H(P))
\]
から，$x,N,p,z,P$ に一様な $0<\lambda_0\le\Lambda_0$ により
\[
 \lambda_0\|P\|\le F(N+P,p,z,x)-F(N,p,z,x)
                   \le\Lambda_0\|P\|\qquad(P\ge0)
\]
を得る．定数の依存も明示できる．\(\lambda I\le H(x)\le\Lambda I\) とし，
\(\|N\|,\|P\|\le\lambda/4\) に制限する．
複素・実座標の対応から
\(\|\mathcal H(Q)\|\le\tfrac12\|Q\|\) であるため，
全ての \(0\le s\le1\) で
\[
 \frac{\lambda}{2}I\le B_s\le2\Lambda I,\qquad
 \frac{\lambda^2}{4\Lambda^2}
 \le\frac{\det B_s}{\det H}
 \le\frac{4\Lambda^2}{\lambda^2}.
\]
さらに \(P\ge0\) に対して
\[
 \frac{\|P\|}{8\Lambda}
 \le\operatorname{tr}_{\CC}(B_s^{-1}\mathcal H(P))
 \le\frac{2\|P\|}{\lambda}.
\]
したがって積分した determinant operator の増分は
\[
 \frac{\lambda^2}{32\Lambda^3}\|P\|
 \le F(N+P,p,z,x)-F(N,p,z,x)
 \le\frac{8\Lambda^2}{\lambda^3}\|P\|.
\]
この \(\lambda,\Lambda\) は固定模型 \(H\) の上下界であり，
actual 解の tensor bounds を先取りしていない．
これは Savin の H2 が用いる実対称方向の operator norm に
関する条件そのものであり，actual $\omega_\nu$ の tensor 上下比較を
入力にしていない．
同じ近傍で $F$ は $C^2$ で，その二階微分は $H$ の一様 $C^2$ bounds によって制御され，
\[
 F(0,0,0,x)=0
\]
を満たす．従って hypothesis H4 も成立する．
H4 は $x$ および zero-order variable $z$ への依存を許している．
ここでは全ての定数が解であるという hypothesis H3 を要求しない．

$v_\nu$ は \eqref{moving:eq:Savin-operator} の滑らかな解，従って viscosity solution である．
Savin\cite[Theorem~1.3]{SavinSmall} を適用すると，小さい $C^0$ norm から
内側 ball 上の一様 $C^{2,\alpha}$ bound を得る．
ここでは同定理の近傍半径 $\delta$ を予め $\lambda$ に比べ十分小さく取り，
結論 $\|v_\nu\|_{C^{2,\alpha}}\le\delta$ から
$H+((v_\nu)_{i\bar j})\ge(\lambda/2)I$ となるようにする．
この段階で \eqref{moving:eq:small-KE-equation} は一様楕円的となるので，
通常の Schauder bootstrap により全固定階の bounds を得る．

率を確認するため，さらに小さい ball 上で
\eqref{moving:eq:small-KE-equation} の logarithm を線形化する：
\[
 A_\nu^{i\bar j}\partial_i\partial_{\bar j}v_\nu-v_\nu=0,
 \qquad
 A_\nu^{i\bar j}
 =\int_0^1\bigl(H+s(v_\nu)_{a\bar b}\bigr)^{-1\,i\bar j}\,ds.
\]
これらの係数は一様楕円性と全固定階の bounds を持つ．
この線形方程式の内点 Schauder estimates により
$\|v_\nu\|_{C^{k+2}(B')}\le C_{k,B'}\|v_\nu\|_{C^0(B)}$
を得る．$\ddc v_\nu=\omega_\nu-h$ から
\eqref{moving:eq:linear-smallness} が従う．
\end{proof}

\begin{theorem}[Dominant triple 領域の actual product-cusp 漸近]
\label{coarse:thm:moving-triple-tensor}
\Cref{ass:Gd} を満たす固定次数 $d\ge5$ の全平面退化を考える．
$a=a_t$, $b=b_t$, $\rho=\rho_t\to\infty$ が
\eqref{moving:eq:rectangle-protection} を満たすとする．
任意の固定 $r>0$ と $k\ge0$ に対して
\begin{equation}\label{moving:eq:actual-tensor-asymptotic}
 \|g_t-h_{\mathrm{prod}}\|_{C^k(D_t(a,b;r),h_{\mathrm{prod}})}
 \le C_{k,r}e^{-\rho_t}
\end{equation}
が十分小さい $t$ に対して成り立つ．
$C^k$ norm は angular universal cover の normalized logarithmic coordinates で計算し，
評価は全 torus orbits について一様である．
有限個の三重点 chart の定数の最大を取ることにより，固定した族の全三重点にも一様にできる．
定数は $d$ と退化族に依存してよい．
特に，この領域では $g_t$ と $h_{\mathrm{prod}}$ の一様 tensor 比較が
独立の仮定なしに成立する．
\end{theorem}

\begin{proof}
任意の固定 $r$ に対し，\cref{coarse:lem:annular-volume-squeeze} を
$r$ より少し大きい固定 inner width に適用する．
Angular universal cover 上で
\[
 w_1=-a^{-1}\log \zeta_1,\qquad
 w_2=-b^{-1}\log \zeta_2
\]
と置けば，$h_{\mathrm{prod}}$ は固定した二つの half-plane metrics
\[
 \frac{|dw_1|^2}{(\Re w_1)^2}
 +\frac{|dw_2|^2}{(\Re w_2)^2}
\]
の積になる．Fixed inner width 上では $\Re w_i$ は正の定数の間にあり，
imaginary translations によりどの位相でも同じ固定 ball を使える．
この ball 上の
$v_t=\log(\omega_t^2/h_{\mathrm{prod}}^2)$ は
$C^0$ norm $O_r(e^{-\rho_t})$ を持つ．
\Cref{coarse:lem:small-volume-tensor} を適用して
\eqref{moving:eq:actual-tensor-asymptotic} を得る．
\end{proof}

\begin{corollary}[Cubical open-square compacta の一様比較]
\label{coarse:cor:open-square-tensor}
同じ固定次数 $d\ge5$ の族で $0<\eta<1/2$ を固定する．
\[
 \frac12L_t^\eta\le a,b\le2L_t^{1-\eta},
 \qquad \rho_t=\frac{\eta}{4}R_t
\]
とすると，\eqref{moving:eq:rectangle-protection} は十分小さい $t$ に対して一様に成り立つ．
従って
\begin{equation}\label{moving:eq:open-square-region}
 E_{t,\eta}
 =\{L_t^\eta\le\ell_1,\ell_2\le L_t^{1-\eta},\quad
      \ell_0=L_t-\ell_1-\ell_2\}
\end{equation}
の全位相を含む近傍で，任意の固定 $k$ に対し
\begin{equation}\label{moving:eq:open-square-rate}
 \|g_t-h_{\mathrm{prod}}\|_{C^k(E_{t,\eta},h_{\mathrm{prod}})}
 \le C_{k,\eta}L_t^{-\eta/4}.
\end{equation}
この領域の raw image は，$\mathcal C_d$ における dominant local label $0$ の正方形の
$(q_1,q_2)\in[\eta,1-\eta]^2$ を $o(1)$ の誤差で実現する．
\end{corollary}

\begin{proof}
\[
 e^{-\rho_t}\min(a,b)\ge\tfrac12L_t^{3\eta/4}\longrightarrow\infty,
 \qquad
 e^{\rho_t}(a+b)\le4L_t^{1-3\eta/4}=o(L_t).
\]
従って定理を全ての中心 $(a,b)$ と位相へ一様に適用できる．
また $q_j=R_t^{-1}\log\ell_j+o(1)$ $(j=1,2)$，
$q_0=1+O(R_t^{-1})$ なので raw image の主張も従う．
\end{proof}

\begin{remark}[閉セルと大域 GH 極限との区別]
\label{moving:rem:scope}
本節で扱うのは二つの独立 depths が同時に $o(L_t)$ となる領域である．
例えば balanced triple $\ell_0=\ell_1=\ell_2=L_t/3$ では
\eqref{moving:eq:rectangle-protection} を保ったまま $\rho_t\to\infty$ とできない．
従って本節の結論だけから全閉 corner の一様比較，source の直径評価，
Lagrangian fibrations の全体での接着，または大域 GH 収束は従わない．
一方，open-square compacta については actual metric の sharp な比較が得られており，
そこに限定した短い torus averaging--Moser 構成には独立の tensor 仮定を要しない．
\end{remark}
